\section{引言}
大语言模型（LLMs）正在快速进步~\cite{gpt4, o1, bc2,omni15}，推动其在医疗领域的更广泛采用~\cite{bcm1,baichuan-m2,MedGemma,Lingshu,Hulu-Med}。因此，业界的期望正在从一次性问答转向端到端的临床决策支持~\cite{MedDM,ArgMed}。这一趋势在 OpenAI 的 GPT-5.2~\cite{openai2025gpt52}、ChatGPT Health~\cite{chatgpt_health} 以及 Claude in Healthcare~\cite{anthropic_healthcare} 等领先系统中体现得尤为明显。然而关键局限仍然存在：尽管在静态、定义明确的基准上得分更高，模型在开放式临床互动中往往难以保持证据约束与不确定性意识，而信息缺失与长时程决策使幻觉更难控制~\cite{hallubug}。因此，我们的目标是超越可靠问答，迈向能够在实践中安全运行的决策支持伙伴。

为弥合被动检索与主动临床支持之间的鸿沟，近期研究通常沿两条部分割裂的范式评估：（i）以事实性为核心的单轮基准；（ii）以流程为导向的多轮问诊模拟。HealthBench~\cite{healthbench} 与 Med-HALT~\cite{ankit2023medhalt} 等基准衡量事实一致性与常见错误模式（包括幻觉）。它们往往显示：在需要复杂、跨约束临床推理的困难病例上性能下降，模型也可能生成缺乏依据的断言。与此同时，OSCE 风格（Objective Structured Clinical Examination）的互动病史采集（IHT）受到越来越多关注，用于评估模型能否主动补全缺失信息并遵循临床流程。Google 的 AMIE~\cite{tao2024amie, elahe2025amie} 等系统在这些规则下展示了较强的模拟沟通质量。

然而，统一这两个维度仍存在关键缺口。现有方法往往将对话交互与临床推理视为相互正交的目标，而不是同一连贯系统的组成部分。知识中心型模型表现为“问诊惰性”，缺乏主动补充缺失证据的能动性；而交互导向模型可能牺牲诊断深度，优先追求流畅性而弱化有原则的鉴别诊断推理。

由于三个技术瓶颈，交互与推理的整合颇具挑战。第一，多样临床任务下的异构训练环境阻碍了稳定的多任务融合。第二，长时程交互放大了强化学习（RL）的归因问题~\cite{guo2025deepseek,bcalign,DBLP:conf/cacre/CaoZ19}：当监督主要来自终局结果时，模型难以识别哪些对话轮次对诊断成功具有因果贡献。第三，提高推理深度常会遇到奖励饱和，性能接近平台期时学习信号减弱——模型为了满足复杂逻辑约束，反而可能增加幻觉~\cite{DBLP:journals/corr/abs-2505-23646,DBLP:journals/corr/abs-2510-22977}。

为解决上述挑战，我们提出 Baichuan-M3——新一代医学 LLM，旨在统一临床问诊与可靠决策。Baichuan-M3 强调与真实临床流程对齐的三项核心能力：（i）主动获取信息；（ii）构建连贯的推理轨迹；（iii）自适应抑制幻觉。

方法上，我们提出由任务特定强化学习（TaskRL）、离线策略蒸馏与多教师在线策略蒸馏（MOPD）组成的三阶段训练框架。该分层设计在融合为单一策略之前先解耦并优化各项能力。为提升长时程问诊表现，我们提出分段流水线强化学习，将复杂任务拆分为多阶段并配置独立奖励信号。同时，我们通过动态量表演化与针对幻觉抑制的 RL 目标进一步强化诊断推理。

Baichuan-M3 在三项权威基准上取得最先进结果：（1）HealthBench-Hard 得分 44.4，超越 GPT-5.2；（2）在 ScanBench（我们的 OSCE 风格基准）的三项维度上均居首位——临床问诊（74.9）、检验检查（72.1）与诊断（74.4）——超过 GPT-5.2-High 与专家基线；（3）在无工具的幻觉评测中表现出更高的事实可靠性。
综上，我们的贡献如下：
\begin{itemize}
    \item \textbf{打通问诊与推理：} 我们提出 Baichuan-M3，旨在将 LLM 从被动信息检索者转变为稳健的决策支持伙伴。通过建模完整的临床决策流程，系统能够主动补全缺失数据，同时保持严谨的诊断逻辑，解决开放式场景中“基于假设的幻觉”这一关键限制。

    \item  \textbf{临床流程对齐的优化：} 我们提出一种模拟专业医学训练的范式，采用分段流水线 RL 使模型行为与临床问诊各阶段（问诊、检验检查、诊断）对齐。该方法结合动态量表演化，确保模型优先学习循证推理，而非单纯的对话流畅或强行逻辑拟合。

    \item \textbf{卓越的实证结果：} 大量实验表明，Baichuan-M3 在事实可靠性与交互流程严谨性上均处于领先水平。它在 HealthBench-Hard 与我们提出的 ScanBench 上刷新 SOTA，在幻觉抑制与诊断准确率方面相较 GPT-5.2 等领先闭源模型与专家基线取得显著提升。
\end{itemize}